\documentclass[11pt]{article}
\usepackage[utf8]{inputenc}
\usepackage[T1]{fontenc}
\usepackage{lmodern}
\usepackage[margin=1in]{geometry}
\usepackage{amsmath,amssymb,amsthm}
\usepackage{enumitem}
\usepackage{booktabs}
\usepackage{xcolor}
\usepackage{tikz}
\usepackage[colorlinks=true,bookmarks=false,linkcolor=blue!50!black,urlcolor=blue!50!black]{hyperref}

\theoremstyle{plain}
\newtheorem{theorem}{Theorem}
\newtheorem{lemma}[theorem]{Lemma}
\newtheorem{claim}[theorem]{Claim}
\newtheorem{corollary}[theorem]{Corollary}
\theoremstyle{definition}
\newtheorem{definition}[theorem]{Definition}
\newtheorem{example}[theorem]{Example}
\theoremstyle{remark}
\newtheorem*{remark}{Remark}

\newcommand{\N}{\mathbb{N}}
\newcommand{\Z}{\mathbb{Z}}
\newcommand{\R}{\mathbb{R}}
\newcommand{\Q}{\mathbb{Q}}
\newcommand{\iv}[1]{[#1]}                 % Iverson bracket
\newcommand{\powerset}{\mathcal{P}}
\DeclareMathOperator{\ord}{ord}
\DeclareMathOperator{\gcdop}{gcd}

\setlength{\parindent}{0pt}
\setlength{\parskip}{0.6em}
\newcommand{\lesson}[2]{\href{../lessons/#1}{#2}}   % link to a lesson image

% Let long inline lists like \{1,2,5,6\} break across lines after a comma.
\mathchardef\mycomma="613B
\mathcode`,="8000
{\catcode`,=\active \gdef,{\mycomma\penalty0\hspace{0pt}}}
\emergencystretch=2em

\title{Van der Waerden's theorem}
\author{Notes from Carlos' lessons}
\date{}
\begin{document}
\maketitle

\textbf{Source lesson:} \lesson{2026-04-23-van-der-waerden.png}{2026-04-23 (van der Waerden)}.

\medskip
\textbf{Question on the board.} If we colour $1,2,\dots,N$ with $c$ colours, can we force a long monochromatic arithmetic progression (AP)? $W(c,L)$ is the least $N$ such that every $c$-colouring of $\{1,\dots,N\}$ contains a monochromatic AP of length $L$.

\begin{theorem}[van der Waerden]
$W(c,L)$ exists for all $c,L$.
\end{theorem}

\section{$W(2,3)=9$}

\paragraph{Lower bound ($W(2,3)>8$).} Colour $1,\dots,8$ as R R B B R R B B. Red is $\{1,2,5,6\}$ and blue is $\{3,4,7,8\}$. No three elements of $\{1,2,5,6\}$ form an AP (check the four triples $\{1,2,5\},\{1,2,6\},\{1,5,6\},\{2,5,6\}$: the gaps are $1,3$ / $1,4$ / $4,1$ / $3,1$, never equal), and $\{3,4,7,8\}$ is the same set shifted by $2$. So there is no monochromatic $3$-AP. (The board's ``Part 1''; verified by brute force.)

\paragraph{Upper bound ($W(2,3)\le9$).} Every $2$-colouring of $1,\dots,9$ has a monochromatic $3$-AP (the board's ``Part 2'').

\begin{proof}[Proof (case analysis)]
Suppose not. By symmetry of the colours let $5$ be red. Each of $\{4,6\},\{3,7\},\{2,8\},\{1,9\}$ then contains a blue element (otherwise there would be a red $3$-AP centred at $5$). Consider $4$ and $6$:
\begin{enumerate}[nosep]
  \item \emph{Both blue.} The APs $(2,4,6)$ and $(4,6,8)$ force $2$ and $8$ to be red; then $(2,5,8)$ is a red AP. Contradiction.
  \item \emph{Both red.} Then $(4,5,6)$ is red. Contradiction.
  \item \emph{Exactly one blue.} The reflection $x\mapsto10-x$ preserves APs and fixes $5$, so assume $4$ is blue and $6$ is red.
  \begin{itemize}[nosep]
    \item $(5,6,7)$: $5,6$ red, so $7$ is blue.
    \item $(1,4,7)$: $4,7$ blue, so $1$ is red.
    \item $(1,5,9)$: $1,5$ red, so $9$ is blue.
    \item $(7,8,9)$: $7,9$ blue, so $8$ is red.
    \item $(2,5,8)$: $5,8$ red, so $2$ is blue.
    \item $(2,3,4)$: $2,4$ blue, so $3$ is red.
    \item Now $1,3,5$ are all red: an AP. Contradiction.
  \end{itemize}
\end{enumerate}
\end{proof}

\section{Why $W(c,L)$ exists}

\paragraph{What the board does.} It sketches the ``columns'' idea: cut $\{1,\dots,N\}$ into blocks of length $m$; a block of length $m$ has one of $c^m$ colourings, so if $N/m\ge W(c^m,L)$, the blocks themselves, viewed as $c^m$ ``super-colours'', contain an AP of $L$ \emph{identical} blocks (``an arithmetic sequence of identical columns''). For $W(2,3)$ the board uses blocks of length $5$ ($2^5=32$ colourings, so any $33$ blocks contain two equal ones; the third block of the AP is the next in the progression). It then tries to turn an AP of length $L$ into one of length $L+1$ and writes a bound involving $W(2,L)$ and $W(2^{W(2,L)},L)$. I could not read that final inequality reliably, and it is stated for $c=2$ only, so I give the standard complete argument instead.

\paragraph{Full proof} (double induction; the ``colour-focusing'' argument as in Graham--Rothschild--Spencer; see the sources below). Induct on $L$: $L=1,2$ are trivial. Assume $W(c,L-1)$ exists for \emph{all} $c$. Fix $c$.

\begin{definition}[Colour-focused APs]
APs $A_i=\{a_i+jd_i:0\le j\le L-2\}$, $i=1,\dots,r$, are \emph{colour-focused at $f$} if each $A_i$ is monochromatic, the colours are pairwise distinct, and $a_i+(L-1)d_i=f$ for all $i$.
\end{definition}

\begin{claim}[$r$]
For each $r\le c$ there is $M(r)$ such that every $c$-colouring of $[1,M(r)]$ contains a monochromatic $L$-AP or $r$ colour-focused $(L-1)$-APs with focus $f\le M(r)$.
\end{claim}

\begin{itemize}[nosep]
  \item \emph{Claim($c$) finishes the proof.} The focus $f$ has one of the $c$ colours, which equals the colour of some $A_i$, and then $A_i\cup\{f\}$ is a monochromatic $L$-AP.
  \item \emph{Claim($1$).} By the induction hypothesis $[1,W(c,L-1)]$ contains a monochromatic $(L-1)$-AP $A_1=\{a_1+jd_1\}$ with $a_1+(L-2)d_1\le W(c,L-1)$, hence $d_1\le W(c,L-1)$, and its focus $a_1+(L-1)d_1\le 2W(c,L-1)=:M(1)$. One colour-focused AP is just a monochromatic $(L-1)$-AP.
  \item \emph{Claim($r-1$) $\Rightarrow$ Claim($r$).} Let $M=M(r-1)$, $m=M$, and $N=W(c^{m},L-1)$. Cut $[1,(2N+1)m]$ into blocks of length $m$ and give each block the ``super-colour'' given by its colouring pattern, one of $c^m$ values. By the induction hypothesis on $L$, among the first $N$ blocks there are blocks $B_0,B_1,\dots,B_{L-2}$ with indices $s,s+t,\dots,s+(L-2)t$ ($t\ge1$) that are \emph{identically coloured}. Inside $B_0$ apply Claim($r-1$) to the relative positions $1,\dots,m$: either there is a monochromatic $L$-AP (done), or there are $r-1$ colour-focused $(L-1)$-APs with relative starts $a_i$, steps $d_i$, and common relative focus $f=a_i+(L-1)d_i\le m$. Since no monochromatic $L$-AP arose, $f$'s colour differs from every $A_i$'s colour.

  Define, in absolute positions, for $i<r$: $A_i'=\{sm+a_i+k(d_i+tm):0\le k\le L-2\}$. Its $k$-th element is the relative position $a_i+kd_i$ inside block $B_k$ (block $B_k$ starts at $(s+kt)m$), and $B_k$ is coloured like $B_0$, so $A_i'$ is monochromatic in the colour of $A_i$. Define $A_r'=\{sm+f+ktm:0\le k\le L-2\}$, the position $f$ in each $B_k$: it is monochromatic in the colour of $f$. The $r$ colours are distinct, and each $A_i'$ has focus $sm+f+(L-1)tm$ (for $i<r$ because $a_i+(L-1)d_i=f$). That focus is at most $(s+(L-1)t+1)m\le(2N+1)m=:M(r)$. So Claim($r$) holds.\hfill$\blacksquare$
\end{itemize}

\section{Aside on the same board: the IMO 1988 Vieta-jumping problem}

The box at the very top of the board is an unrelated problem. I read it as the numerator and denominator written with the roles swapped, but it is surely this one: \emph{if $a,b$ are positive integers and $k=\dfrac{a^2+b^2}{ab+1}$ is an integer, then $k$ is a perfect square} (IMO 1988, Problem 6).

\begin{proof}[Proof (Vieta jumping)]
Suppose $k$ is not a square. Among all positive-integer solutions of $a^2-kab+b^2=k$ take one with $a+b$ minimal, and $a\ge b$. As a quadratic in $a$ it has roots $a$ and $a'=kb-a=\dfrac{b^2-k}{a}$. The number $a'$ is an integer. If $a'=0$ then $b^2=k$, a square. If $a'\le-1$ then $a'^2-ka'b+b^2\ge1+k+1>k$, contradicting the equation. So $a'\ge1$, and $a'=\frac{b^2-k}{a}<\frac{b^2}{a}\le b\le a$. Then $(a',b)$ is a solution with a smaller sum. Contradiction.
\end{proof}

\section*{Sources}
\begin{itemize}[nosep]
  \item \href{https://en.wikipedia.org/wiki/Van_der_Waerden\%27s_theorem}{Van der Waerden's theorem (Wikipedia)}: the value $W(2,3)=9$ and the double-induction outline.
  \item \href{https://opentext.uleth.ca/Ramsey/sec_proof_vdW.html}{Proof of van der Waerden's theorem (opentext.uleth.ca)} and Blondal--Juni\'c, \emph{Proof of van der Waerden's Theorem in Nine Figures}: the colour-focusing proof.
  \item The $W(2,3)$ case analysis above is my own and was checked by brute force: the extremal colouring of $1,\dots,8$ is R R B B R R B B (or its colour swap).
\end{itemize}

\end{document}
